\section{Recollections}\label{recollections}


Associated to the square 
\[
\begin{tikzcd}
\BP\langle 1\rangle \ar[r] \ar[d] & H\mathbb{Z}_{(p)} \ar[d] \\ 
k(1) \ar[r] & H\mathbb{F}_p
\end{tikzcd}
\]
there is a square of Bockstein spectral sequences \[
\begin{tikzcd}
\THH_*((\BP\langle n \rangle ;\mathbb{F}_p)[v_0,v_1]\arrow[Rightarrow]{r} \arrow[Rightarrow]{d} &  \THH_*(\BP\langle n\rangle ;\mathbb{Z}_{(p)})[v_1]\arrow[Rightarrow]{d}\\ 
\THH_*(\BP\langle n\rangle ;k(1))[v_0]\arrow[Rightarrow]{r}  &  \THH_*(\BP\langle n \rangle ;\BP\langle 1\rangle) 
\end{tikzcd}
\]
for any integer $n\ge 1$. The first Bockstein spectral sequence in each composite of spectral sequences above was computed in~\cite{MS93},~\cite{AR05} and~\cite{AHL10} for $n=1$ and~\cite{ACH24} for $n=2$. We recall some of these results below. 

% First, let 
% \[ 
% r(s,n) = \begin{cases} 
% 0 \text{ for }s\le 0, \text{ and } \\ 
% p^{s}+r(s-(n+1),n) \text{ for }s>0\,. 
% \end{cases}
% \]

% \begin{theorem}[{\cite{MS93}},{\cite{AR05}}]
% The $\mathbb{F}_p[v_1]$-module $\THH_*(\BP\langle 1\rangle; k(1))$ is generated by $1$, $x_{n,m}=\lambda_n\mu^{p^{n-1}m}$, $x'_{n,m}=\lambda_n\lambda_{n+1}\mu^{p^{n-1}m}$ for $n\ge 1$ and $m\ge 0$, $m\ne p-1\mod p$. 
% The relations are generated by 
% \[ 
% v_1^{r(n,1)}x_{n,m}=v_1^{r(n,1)}x'_{n,m}=0.
% \]
% \end{theorem}

\begin{theorem}[{\cite[\S~3.2]{AHL10}}]
The homotopy of $\THH(\BP\langle 1\rangle;\mathbb{Z}_{(p)})$ is a copy of $\mathbb{Z}_{(p)}$ generated by $\lambda_1$ plus torsion. The torsion is generated as a $\mathbb{Z}_{(p)}$-module by the
elements $a_i$ and $b_i$ where $a_i=\mu_2^{i-1}\lambda_2$ and $b_i=\mu_2^{i-1}\lambda_1\lambda_2$ for $i\ge 1$ and they both have order $p^{k+1}$, where $k= \nu_{p}(i)$, the
$p$-adic valuation of $i$. Here $|a_i|=2p^2i-1$ and $|b_i|=2p^2i+2p-2$
\end{theorem}

\begin{theorem}[{\cite{ACH24}*{Theorem~3.8}}]\label{thm:BP2-coeff}
Suppose the error term \cite{ACH24}*{(3--7)} vanishes at $p>3$. There is an isomorphism 
\[ 
\THH_*(\BP\langle 2\rangle ;\mathbb{Z}_{(p)})=\mathbb{F}_p\langle \lambda_1 ,\lambda_2\rangle \otimes \left ( \mathbb{Z}_{(p)}\oplus T_0^2 \right )
\]
where 
\[ T_0^2 = \oplus_{s\ge 1} \mathbb{Z}/p^s\otimes \mathbb{Z}_{(p)}[\mu_3^{p^s}]\otimes \mathbb{Z}_{(p)}\{\lambda_{s+2}\mu_3^{jp^{s-1}} \mid 0\le j\le p-2 \} \,. 
\]
% \item There is an isomorphism
% \[
% \THH_*(\BP\langle 2\rangle ;k(1))= \mathbb{F}_p\langle \lambda_1 \rangle \otimes \left ( \mathbb{F}_p[v_1]\oplus T_1^2 \right ) 
% \]
% where 
% \[ 
%   T_1^2=\bigoplus_{s\ge 1}\mathbb{F}_p[v_1]/(v_1^{r(s,1)})\otimes \mathbb{F}_p[\mu_3^{p^s}]\otimes \mathbb{F}_p\langle \lambda_{s+2}\rangle \otimes \mathbb{F}_p\{ \lambda_{s+1}\mu_3^{jp^{s-1}} \mid 0\le j\le p-2 \} \,.
% \]
% \end{enumerate}
Here $|\mu_3|=2p^3$ and $|\lambda_{s}|=2p^{s}-1$. 
\end{theorem}
Note that the error term \cite{ACH24}*{(3--7)} vanishes for $\mathbb{E}_\infty$ algebras so the theorem holds unconditionally for $\BP\langle 2\rangle=\taf^{D}$ from~\cite{HL10}. 

The map 
\[ 
  \THH(\BP\langle 2\rangle ; \mathbb{F}_p)\to  \THH(\BP\langle 1\rangle ; \mathbb{F}_p) 
\]
sends $\lambda_s$ to $\lambda_s$ for $s=1,2$ and $\mu_3$ to $\mu_2^p$. For consistency, let $a_{pi+1}=\lambda_1\mu_3^{i}$ and $b_{pi+1}=\lambda_1\lambda_2\mu_3^{i}$ 

\begin{corollary}
  The map 
  \[ 
    \THH_*(\BP\langle 2\rangle ;\mathbb{Z}_{(p)})\longrightarrow \THH_*(\BP\langle 1\rangle ;\mathbb{Z}_{(p)})
  \]
  sends $\lambda_1$ to $\lambda_1$, and it sends $\lambda_1\mu_3^i$ to $a_{pi+1}$ and $\lambda_1\lambda_2\mu_3^i$ to $b_{pi+1}$. 

  The map 
  \[ 
    \THH_*(\BP\langle 2\rangle ;k(1))\longrightarrow \THH_*(\BP\langle 1\rangle ;k(1))
  \]
  sends $x_{n,m}$ to $x_{n,m}$ and $x'_{n,m}$ to $x'_{n,m}$ whenever this makes sense.
\end{corollary}



We recall the structure of $\THH_*(\BP \langle 1 \rangle) = \THH_*(\ell)$.

\begin{theorem}[{\cite[Sections 6.2--6.3]{AHL10}}]\label{thm:AHL}
  $\THH_*(\BP \langle 1 \rangle)$ is a quotient of the $\Z_{(p)}[v_1]$-module
  \begin{multline}
    \Z_{(p)}[v_1]\{1,\, \lambda_1,\, v_0^n a_{p^n},\,n\geq 0\} 
    \oplus \Z_{(p)}[v_1]\{v_0^h b_{\alpha p^n},\, n\geq 0,\, \alpha\geq 1,\, p\nmid \alpha,\, h \geq 0\}
  \end{multline}
  by the relations in the non-torsion part:
  \begin{itemize}
  \item $p\cdot a_1 = v_1^p \lambda_1$,
  \item $p\cdot v_0^n a_{p^n} = v_1^{p^{n+1}}v_0^{n-1} a_{p^{n-1}}$ for any $n\geq 1$,
  \end{itemize}
  and the relations in the torsion part:
  \begin{itemize}
  \item $v_0^{h} b_{\alpha p^n} = 0$ for any $\alpha \geq 1$ and $n \geq 0$, $\alpha$ not divisible by $p$, and $h \geq n + 1$,
  \item $v_1^{p^{n-h+1}+p^{n-h}+\dots+p }\cdot v_0^h b_{\alpha p^n} = 0$ for any $\alpha \geq 1$ and $n \geq 0$, $\alpha$ not divisible by $p$ and  $0 \leq h \leq n$,
  \item $p\cdot b_{(\beta p+p-1)p^n} = v_0 b_{(\beta p+p-1)p^n}+v_1^{p^{n+2}}v_0^{\nu_p(\beta)} b_{\beta p^{n+1}}$ for any $\beta \geq 1$ and $n \geq 0$. Here $\nu_p$ denotes $p$-adic valuation.
  \item $p\cdot v_0^h b_{\alpha p^n} = v_0^{h+1} b_{\alpha p^n}$ for any $\alpha \geq 1$, $n \geq 0$, $\alpha$ not divisible by $p$, and any $1 \leq h \leq n$, or $h=0$ not in the previous case.
  \end{itemize}
Here $|\lambda_1|=2p-1$, $|v_0^na_{p^n}|=2p^{n+2}-1$, $|v_1|=2p-2$ and $|v_0^hb_{\alpha p^n}|=2p^{n+2}\alpha+2p-2$.
\end{theorem}

We also recall the computation of the rationalization of topological Hochschild homology of truncated Brown--Peterson spectra. 

\begin{proposition}[{\cite{ACH24}*{Proposition~3.7}}]\label{prop:rational}
Let $0\le m\le n$. There is a preferred isomorphism 
\[ 
\mathbb{Q}[v_1,\cdots ,v_m]\langle\sigma v_1,\sigma v_2,\cdots ,\sigma v_n) \cong \THH_*(\BP\langle n\rangle ; \BP\langle m\rangle)_{\mathbb{Q}} 
\]
where we write $X_{\mathbb{Q}}$ for the Bousfield localization at the Eilenberg--MacLane $H\mathbb{Q}$. 
\end{proposition}
